%-------------------------
% Resume in Latex
% Author: Garv Arora
% Based off of: https://github.com/sb2nov/resume
% License: MIT
%------------------------

\documentclass[letterpaper,9pt]{extarticle}

\usepackage{latexsym}
\usepackage[empty]{fullpage}
\usepackage{titlesec}
\usepackage{marvosym}
\usepackage[usenames,dvipsnames]{color}
\usepackage{verbatim}
\usepackage{enumitem}
\usepackage[hidelinks]{hyperref}
\usepackage{fancyhdr}
\usepackage[english]{babel}
\usepackage{tabularx}
\usepackage[top=0.4in, bottom=0.4in, left=0.5in, right=0.5in]{geometry}
\input{glyphtounicode}

\pagestyle{fancy}
\fancyhf{}
\fancyfoot{}
\renewcommand{\headrulewidth}{0pt}
\renewcommand{\footrulewidth}{0pt}

\urlstyle{same}
\raggedbottom
\raggedright
\setlength{\tabcolsep}{0in}

\titleformat{\section}{
  \vspace{-4pt}\scshape\raggedright\large
}{}{0em}{}[\color{black}\titlerule\vspace{-4pt}]

\pdfgentounicode=1

%-------------------------
% Custom commands

\newcommand{\resumeItem}[1]{
  \item\small{#1 \vspace{-1pt}}
}

\newcommand{\resumeSubheading}[4]{
  \item
    \begin{tabular*}{0.97\textwidth}[t]{l@{\extracolsep{\fill}}r}
      \textbf{#1} & \small #2 \\
      \textit{\small#3} & \textit{\small #4} \\
    \end{tabular*}\vspace{-5pt}
}

\newcommand{\resumeSubRole}[2]{
  \vspace{3pt}
  \begin{tabular*}{0.97\textwidth}[t]{l@{\extracolsep{\fill}}r}
    \textit{\small#1} & \textit{\small #2} \\
  \end{tabular*}\vspace{-5pt}
}

\newcommand{\resumeProjectHeading}[2]{
  \item
    \begin{tabular*}{0.97\textwidth}{l@{\extracolsep{\fill}}r}
      \small#1 & \small #2 \\
    \end{tabular*}\vspace{-5pt}
}

\renewcommand\labelitemii{$\vcenter{\hbox{\tiny$\bullet$}}$}

\newcommand{\resumeSubHeadingListStart}{
  \begin{itemize}[leftmargin=0.15in, label={}, itemsep=5pt, parsep=0pt, topsep=3pt]
}
\newcommand{\resumeSubHeadingListEnd}{\end{itemize}}

\newcommand{\resumeItemListStart}{
  \begin{itemize}[leftmargin=0.2in, itemsep=1pt, parsep=0pt, topsep=3pt]
}
\newcommand{\resumeItemListEnd}{\end{itemize}\vspace{1pt}}

%-------------------------------------------
\begin{document}

%----------HEADING----------
\begin{center}
    \textbf{\Huge \scshape Garv Arora} \\ \vspace{2pt}
    \small\textit{Robotics \& Embedded Systems Engineer $|$ Autonomous Systems $|$ ADAS} \\ \vspace{2pt}
    \small Vellore, India $|$
    +91 88008-12254 $|$
    \href{mailto:garvarora0205@gmail.com}{garvarora0205@gmail.com} $|$
    \href{https://linkedin.com/in/gaminization}{linkedin.com/in/gaminization} $|$
    \href{https://github.com/gaminization}{github.com/gaminization}
\end{center}

%-----------SUMMARY-----------
\section{Summary}
\vspace{1pt}
\small{R\&D Lead at SEDS VIT, directing Team Ardra (ISDC) and Team Vyadh (IRC/IRDC) --- 50+ engineers across autonomous drone and planetary rover programmes under SEDS India. Ships production ROS2 navigation stacks, CUDA-accelerated perception systems, mmWave sensor pipelines, and ARM bare-metal RTOS. Placed 12th (IRC~2025) and 16th (IRC~2026) globally among 100+ international teams. Open to roles in robotics software, embedded systems, ADAS, or autonomous vehicle engineering.}

%-----------EDUCATION-----------
\section{Education}
  \resumeSubHeadingListStart
    \resumeSubheading
      {Vellore Institute of Technology $|$ Vellore, India}{CGPA: 8.37 / 10.0}
      {Bachelor of Technology -- Computer Science \& Engineering (Specialization: IoT)}{August 2023 -- July 2027}
    \resumeSubheading
      {Gems International School $|$ Gurugram, India}{Class 10: 89.3\% $|$ Class 12: 76.4\%}
      {Central Board of Secondary Education (CBSE)}{November 2016 -- April 2023}
  \resumeSubHeadingListEnd

%-----------TECHNICAL SKILLS-----------
\section{Technical Skills}
\begin{itemize}[leftmargin=0.15in, label={}, itemsep=1pt, parsep=0pt, topsep=1pt]
  \small{\item{
    \textbf{Languages}{: C, C++, Python, Java, JavaScript, SQL, Bash, Verilog, C\#, HTML/CSS} \\
    \textbf{Robotics \& Autonomous Systems}{: ROS2, Nav2, MoveIt2, RTAB-Map, Micro-ROS, SLAM, EKF/UKF, PID Control, Point Cloud Processing, mmWave Radar} \\
    \textbf{Perception \& ML}{: OpenCV, YOLOv8, MediaPipe, Intel RealSense SDK, PyTorch, TensorFlow, scikit-learn, CUDA, HuggingFace Transformers} \\
    \textbf{Embedded \& Infrastructure}{: ARM Cortex-M, ESP32, I2C/SPI/UART, FreeRTOS, Docker, Git, GitLab CI/CD, Linux, AWS, PyQt, MongoDB}
  }}
\end{itemize}

%-----------WORK EXPERIENCE-----------
\section{Work Experience}
  \resumeSubHeadingListStart

    \resumeSubheading
      {SEDS Projects VIT $|$ SEDS India}{}
      {Research \& Development Lead}{January 2026 -- Present}
      \resumeItemListStart
        \resumeItem{Promoted to R\&D Lead overseeing 2 competition teams (50+ engineers) --- Team Ardra (ISDC) and Team Vyadh (IRC/IRDC) --- under SEDS India, the national apex body for space engineering students.}
        \resumeItem{Guided 3 sub-team leads through navigation architecture reviews and hardware integration debugging, resolving 10+ critical pre-competition issues across both teams.}
      \resumeItemListEnd
      \resumeSubRole{Autonomous Systems Developer -- Team Vyadh (Martian Rover)}{April 2024 -- January 2026}
      \resumeItemListStart
        \resumeItem{Architected a full ROS2 navigation stack (Nav2, RTAB-Map Visual SLAM, RealSense D455) generating real-time 3D elevation maps, cutting manual teleoperation time by 70\% in competition.}
        \resumeItem{Slashed CV pipeline latency by $\sim$60\% and raised throughput 4x via CUDA, YOLOv8, ArUco/AprilTag detection, and Linux V4L2 tuning across 4 simultaneous camera streams.}
        \resumeItem{Flashed Micro-ROS firmware to ESP32 microcontrollers for distributed data acquisition; fused IMU and encoder data with an EKF, cutting odometry drift by $\sim$40\% and enabling sub-10 cm autonomous positioning.}
        \resumeItem{Shipped a PyQt5 autonomous mission dashboard with real-time telemetry visualization, parameter tuning, and FSM-based mission sequencing for autonomous competition runs.}
      \resumeItemListEnd

    \resumeSubheading
      {Wissen Baum Engineering Solutions}{May 2025 -- July 2025}
      {Software Automation Intern}{Pune, India}
      \resumeItemListStart
        \resumeItem{Authored a BDD test suite (Python, Behave, pytest) covering 40+ scenarios, raising regression coverage from 0\% to 100\% on 3 core product flows in 6 weeks.}
        \resumeItem{Accelerated GitLab CI/CD by parallelizing 3 stages (Playwright + Cypress), trimming run time $\sim$35\% (18 min to 12 min); replaced manual QA with a pandas/NumPy script processing 10,000+ rows per run.}
        \resumeItem{Enforced code quality across the shared repository by configuring pre-commit Git hooks with linting rules, blocking 100\% of non-compliant pushes before merge.}
      \resumeItemListEnd

  \resumeSubHeadingListEnd

%-----------PROJECTS-----------
\section{Projects}
  \resumeSubHeadingListStart

    \resumeProjectHeading
      {\textbf{mmWave Earthquake Survivor Detection Robot} $|$ \emph{ROS2, Python, Hi-Link LD2410C, ESP32}}{2026}
      \resumeItemListStart
        \resumeItem{Integrated an LD2410C mmWave radar with ROS2 to detect survivor breathing signatures through rubble; fused radar presence data with IMU odometry for autonomous debris-field navigation, marking locations with $<$15 cm accuracy.}
        \resumeItem{Streamed radar frames from the LD2410C to an ESP32 over UART, published processed detections as ROS2 topics, and visualised survivor coordinates in real time on a 2D occupancy map.}
      \resumeItemListEnd

    \resumeProjectHeading
      {\textbf{Gesture-Controlled 5-DOF Robotic Arm} $|$ \emph{C++, Arduino, MPU6050, Flex Sensors}}{2026}
      \resumeItemListStart
        \resumeItem{Mapped MPU6050 wrist orientation and 2 flex sensor readings to 5-DOF servo commands with $<$50 ms end-to-end latency; trained a gesture classifier on 200+ samples, achieving 94\% accuracy across 8 gestures.}
        \resumeItem{Sent gesture commands over HC-05 Bluetooth to the arm controller, decoupling the glove from the manipulator and enabling wireless operation up to 10 m.}
      \resumeItemListEnd

    \resumeProjectHeading
      {\textbf{HayaiOS -- Bare-Metal RTOS} $|$ \emph{C, ARM Assembly, ARM Cortex-M4}}{2026}
      \resumeItemListStart
        \resumeItem{Implemented a preemptive scheduler (1 ms SysTick), context switching in $<$50 assembly instructions, mutex/semaphore primitives, and a register-level HAL on ARM Cortex-M4.}
        \resumeItem{Stress-tested the scheduler under 8 concurrent tasks, confirming zero priority inversion and $<$2 $\mu$s context-switch latency across 10,000 scheduling cycles.}
      \resumeItemListEnd

  \resumeSubHeadingListEnd

%-----------CERTIFICATIONS & ACHIEVEMENTS-----------
\section{Certifications \& Achievements}
\begin{itemize}[leftmargin=0.15in, label={}, itemsep=1pt, parsep=0pt, topsep=1pt]
  \small{\item{
    \textbf{Oracle Cloud Infrastructure 2025:}
    \href{https://catalog-education.oracle.com/pls/certview/sharebadge?id=EF9A070356E9FB420B39ED42ACCD6CC9326E5653AE07F025C334A2F0238D4652}{Generative AI Professional}
    $|$
    \href{https://catalog-education.oracle.com/pls/certview/sharebadge?id=C3199CDA4E4B76D94AEAE70CFCA292011795FEFB3BCD93A18A1454E57C60B0AA}{Data Science Professional} \\
    \textbf{IRC 2025}{: 12th globally (100+ teams)} $|$
    \textbf{IRC 2026}{: 16th globally} $|$
    \textbf{IRDC 2025}{: Top 5 Finalist}
  }}
\end{itemize}

%-------------------------------------------
\end{document}
